%Os pacotes abaixo adicionam comandos LaTeX que facilitam a escrita e formatação. Caso tenha alguma dúvida sobre a funcionalidade, pesquise o nome do pacote no link: https://www.overleaf.com/learn/latex/Understanding_packages_and_class_files

\usepackage{graphicx}  % Inclui arquivos de figuras
\usepackage{subfigure}
\usepackage{multirow}

\linespread{1.1}
\usepackage{fancyhdr}
\usepackage{longtable}
\usepackage{parskip}
\usepackage[brazilian]{babel}
\usepackage[T1]{fontenc}
\usepackage{float}
\usepackage{dcolumn}
\usepackage{booktabs}   
\usepackage{bm}        
\usepackage{amsfonts}  % Fontes de matemática
\usepackage{amsmath}   % Funções matemáticas
\usepackage{amssymb}   % Símbolos matemáticos
\usepackage{lipsum}

\setlength{\parindent}{10pt}
\newcommand{\ccite}[1]{[\citeauthor{#1}, Ref. \citenum{#1}]}